%Dokumentenklasse
\documentclass[
	a4paper,										%Format
	12pt,											%Schriftgröße
	twoside,										%Zweiseitig
	headsepline,								%horizontale Linie oben
	footsepline,									%horizontale Linie unten
	captions=tableheading,					%Überschrift über Tabelle (nicht unter)
]{scrartcl}

\renewcommand{\thesubsection}{\thesection.\arabic{subsection}.}
\renewcommand{\thesubsubsection}{\thesection.\arabic{subsection}.\arabic{subsubsection}. \hspace{2pt}}
\renewcommand{\theparagraph}{\thesection.\arabic{subsection}.\arabic{subsubsection}.\arabic{paragraph}.}

%Pakete:
\usepackage[exponent-product = \cdot, output-decimal-marker = {,}]{siunitx}
\usepackage{textalpha}
\usepackage{floatrow}
\floatsetup[table]{style=plaintop}
\floatsetup[longtable]{LTcapwidth=table}
%\usepackage{overcite}
\usepackage{siunitx}
\usepackage[ngerman]{babel} 			%Neue deutsche Rechtschreibung
\usepackage{grffile}
\usepackage{graphicx} 					%Bilder einbinden
\usepackage{xcolor}						%Farben für Schrift und Plots
\usepackage{pgfplots}
\pgfplotsset{compat=1.18}					%Plots und Tabellen aus Datensätzen.
\usepackage{amsmath}						%Mathe
\usepackage{amssymb}					%Mathesymbole
\usepackage{longtable}					%Seitenübergreifende Tabellen zulassen
\usepackage{float}
\usepackage{tabularx}
\usepackage{geometry}					%Seitenränder
\usepackage{ae}								%skalierbare Schriften (im PDF-File)
\usepackage[activate]{microtype}		%optischer Randausgleich
\usepackage{setspace}					%Zeilenabstand: \onehalfspacing
\usepackage[utf8]{inputenc}				%Zeichenkodierung für Umlaute
\usepackage{scrlayer-scrpage}					%für Dokumentformatierung mit Pagestyle
\usepackage[version=4]{mhchem}						%Chemische Formeln (mit \ce )
\usepackage{chemfig}
\usepackage{siunitx}
\usepackage{caption}										%Überschriften formatieren
\captionsetup{labelfont={bf,it},tablename=Tab.,figurename=Abb.}				%Überschriften formatieren
\captionsetup[table]{singlelinecheck=false,justification=raggedright}
\captionsetup[figure]{justification=raggedright}				

%Abbildungs und Tabellennummerierung anpassen
\renewcommand*\figureformat {\figurename~\arabic{section}.\arabic{subsection}.\thefigure.}	%Punkt nach Abb. #
\renewcommand*\tableformat {\tablename~\arabic{section}.\arabic{subsection}.\thetable:}		%Doppelpunkt nach Tab. #
\renewcommand*\captionformat {~}								%Nur Abstand, kein Zeichen nach Tab. #: bzw Abb. #. 

%Literatur als Subsubsection definieren und von Literaturverzeichnis in Literatur umbenennen
\makeatletter
\renewcommand*\bib@heading{\subsubsection{\bibname}			
\@mkboth{\subsubsectionmarkformat \bibname}{\subsubsectionmarkformat \bibname}}
\makeatother
\addto\captionsngerman{\renewcommand*\bibname{Literatur}} 

\usepackage{prettyref}						%Nummerierung der Tabellen und Abbildungen im Text an Tab. x.x.# anpassen
%%% Für Abbildungen %%%
\newrefformat{fig}{Abb.~\arabic{section}.\arabic{subsection}.\ref{#1}}
%%% Für Tabellen %%%
\newrefformat{tab}{Tab.~\arabic{section}.\arabic{subsection}.\ref{#1}}

%Links im Dokument:
\usepackage[									%als letztes Paktet einbinden!
	linkcolor=black,							%Farbe der Links
	citecolor=black,							%Farbe der Link-Ziele
	urlcolor=black,								%Farbe von URLs
	colorlinks=true,							%Links nicht als Kasten, sondern direkt als Text anzeigen
]{hyperref}										%Links im PDF-File, unterstützung von URLs mit /url{}

%Dokumentformatierungen
\setkomafont{sectioning}{\bfseries} 	%Serifenschrift in Überschriften
\setkomafont{subsection}{\LARGE}
\setkomafont{subsubsection}{\large}
\pagenumbering{arabic}					%arabische Seitenzahlen
\geometry{left=3cm, right=2cm, top=2.5cm, bottom=2cm} 	%Seitenränder
\parindent0pt 									%Absätze einrücken um 0Pt
\parskip 10pt									%Abstand nach Absatz
\onehalfspacing								%Zeilenabstand
\pagestyle{scrheadings}					%Kopf- und Fussleistenformatierung
\headsep0.4cm								%Abstand Kopfleiste-Text
\footskip0.8cm									%Abstand Fussleiste-Text
\urlstyle{same}								%URLs normal formatieren, nicht als Typewriter

%Befehle definieren

%%%%% create chemequation environment

\newcounter{chemequation}
\renewcommand*\thechemequation{(\Roman{chemequation})}
\newenvironment{chemequation}{%
	\stepcounter{chemequation}%
	\begin{equation}}%
{\tag*{\thechemequation}%
	\end{equation}}

%%%%% create chemtag to iterate and ref chemequations in an align environment

\newcommand*\chemtag{%
	\stepcounter{chemequation}%
	\tag*{\thechemequation}}

%%%%% to create a chemicalequation float

\usepackage{float}
\newfloat{chemicalequation}{htbp}{cex}%[chapter]
\floatname{chemicalequation}{}

\newcommand{\zitat}[1]{\textsuperscript{\cite{#1}}}

\newcommand{\bild}[4]{					%Bilder mit Beschriftung einfügen
	\begin{figure}[#4]						%#4: Platzierung (htbp)
		\centering								%Bilder mittig einfügen
		\captionbox								%Beschriftungsumgebung
			{#2\label{fig: #1}}				%#2: Beschriftung
			{\includegraphics[#3]{#1}}	%#1: Dateiname, #3: Bildeigenschaften
	\end{figure}} 

\newcommand{\formel}[2]{				%Formel einfügen und Referenz erzeugen
	\begin{equation}
		\label{eq: #1}							%#1: Formelname
		#2											%Formelinhalt
	\end{equation}}

\newcommand{\cformel}[2][]{				%Chemische Formel einügen und Referenz erzeugen
	\begin{chemequation}
		\label{ceq: #1}					%#1: Formelname (optional)
		\ce{#2}							%#2: Formelinhalt
	\end{chemequation}\\}

	% 2. Ein neuer, eigener Befehl NUR für Formeln mit Zusatz auf der rechten Seite:
\newcommand{\cformelpot}[3]{%
    \begin{chemequation}
        \label{ceq: #1}
        \ce{#2} \hspace{2cm} #3
    \end{chemequation}%
}

\newcommand{\ox}[2]{
	\overset{\text{#2}}{\ce{#1}}}
	
\newcommand{\tabelle}[6]{				%Tabelle mit Beschriftung einfügen
\begin{table}[#3]							%#3: Platzierung (htbp)
\centering										%Tabelle mittig einfügen
\captionbox										%Beschriftungsumgebung
{#2												%#2: Tabellenbeschriftung
\label{tab: #1}}								%#1: Tabellenname
{\begin{tabular}{#4}						%#4: Anzahl und Ausrichtung Spalten (l c r)
\hline 											%horizontale Linie
#5													%#5: Obere Zeile
\hline
#6													%#6: Datenzeilen
\hline 
\end{tabular}}
\end{table}}

\usepackage{array,ragged2e}
\newcolumntype{P}[1]{>{\RaggedRight\arraybackslash}p{#1}} %Verhindert ungewollte Lücken in Tabellen

\newcommand{\langetabelle}[6]{  %Seitenübergreifende Tabelle, Kopfzeile wird 
	\begin{centering}           %auf jeder neuen Seite ebenfalls angezeigt
		{\begin{longtable}[#3]{#4}
				\caption{#2\label{tab: #1}}\\				
				\hline
				#5
				\hline
				\endfirsthead
                \caption{#2\label{tab: #1}}\\
				\hline
				#5
				\hline
				\endhead
				#6												
				\hline 
		\end{longtable}}
\end{centering}}

%Plots erstellen&&&&&&&&

%#1: Referenz
%#2: Name
%#3: x label
%#4: y label
%#5: x-werte,y-werte (Spaltenweise)
%#6: Quelle (.txt Datei)
%#7: "Nummerierung" (manuell) <.<

\newcommand{\plot}[6]{
	\begin{figure}[H]	
		\centering
		\begin{tikzpicture}
	
		\begin{axis}[
		title style={at={(0.5,0)},
		anchor=north,
		yshift=-40},
		xlabel={#3},
		ylabel={#4},
		legend pos=south east,
		legend entries={},
		]
		\addplot[color=black,mark=*,only marks] table [#5] {#6};
		\label{plot: plot1}
		\end{axis}
		\end{tikzpicture}
		\caption{#2}
		\label{fig: #1}
	\end{figure}
}

\newcommand{\molL}[1]{\ensuremath{\mathrm{\SI{#1}{\frac{mol}{L}}}}} % \molL{1} => \SI{1}{\frac{mol}{L}}
\newcommand{\mL}[1]{\ensuremath{\mathrm{\SI{#1}{mL}}}} % \mL{1} => \SI{1}{mL}

\newcommand{\Rn}[1]{\romannumeral#1\relax}                   %\Rn{1} = i
\newcommand{\RN}[1]{\uppercase\expandafter{\romannumeral#1}} %\RN{1} = I

%Dokument:%%%%%%%%%%%%%%%%%%%%%%%%%%%%%%%%%%%%%%%%%%%%%%%%%%%%%%%%%%%%%%%%%%%%%%%%%%%%
\begin{document}
	
	%Titelseite
	\begin{titlepage}
		\begin{center}
			\hrule													%waagerechte Linie
			\vspace*{5cm}
			\Huge {\textbf {Versuchsreihe 6}}\\						%Titel Versuchsreihe
			\vspace*{0.5cm}
			%		\large {\bf {Untertitel}}\\						%Untertitel
			\vspace*{1cm}
			\normalsize
			Protokoll	\\									
			Praktikum für Nanowissenschaftler:innen\\
			im Studiengang Nanowissenschaften\\
			\vspace*{0.5cm}
			\Large {\textbf{Elif Öner, Friederike Penk, Maxim Popow, Michael Quint}}		\\						%Name
			\vspace*{0.5cm}
					\normalsize{Plätze 7a, 8a, 8b, 9a}								%Platznummer
			\vfill
		\end{center}
		\begin{tabbing}
			\hspace{3cm}\=\kill
			Assistenten: \>Maximilian Ruffer \\						% Name Assistent
			Version: \> 1 \\									%Version
			Datum: \> 17.11.2026 \\								%Datum
		\end{tabbing}
		\hrule														%Waagerechte Linie
\end{titlepage}

\setcounter{page}{1}
\setcounter{section}{8}								%Versuchstag #.x.x bzw Versuchsreihe
\setcounter{subsection}{0}							%Versuchsnummer x.#.x startet bei 0 für 1 
\setcounter{secnumdepth}{4}
	
\subsection{Photometrie}

\subsubsection{Einleitung und Aufgabenstellung}

\subsubsection{Theoretischer Hintergrund}

Die Photometrie ist eine Analysemethode in der Chemie,
mit welcher die Konzentration einer in Lösung befindenden lichtabsorbierenden Substanz ermittelt werden kann.
Dabei werden die Intensitäten der Lichtstrahlen vor und nach dem Durchgang einer Lösung miteinander verglichen.\zitat{kp}

Das Lambert-Beersche Gesetz beschreibt diesen Zusammenhang mathematisch.
Dabei ist $c$ die Konzentration der Lösung, $I$ die austretende Intensität, $I_0$ die eintretende Intensität,
$l$ die optische Weglänge und $\epsilon^*$ der natürliche molare Absorptionskoeffizient.
\formel{a1}{I = I_{0} \cdot \mathrm{e}^{-\epsilon^* \cdot c \cdot l}}
\formel{a2}{\ln(\frac{I}{I_{0}}) = -\epsilon^* \cdot c \cdot l}

Da es allerdings gängig ist, den dekadischen Logarithmus zu verwenden, wird der dekadische molare Absorptionskoeffizient $\epsilon$ eingeführt:
\formel{zusatz}{\epsilon^* = \epsilon \cdot \ln(10)}
\formel{a3}{\lg(\frac{I}{I_{0}}) = -\epsilon \cdot c \cdot l}

Der logarithmische Term wird als die Extintion bzw. als die Absorbanz $A$ definiert,
welche als Maß für das absorbierte (und ebenfalls auch gebeugte und gestreute) Licht fungiert.
\formel{a4}{A = \lg(\frac{I_{0}}{I}) = \epsilon \cdot c \cdot l}

Jede lichtabsorbierende Substanz besitzt dementsprechend einen molaren Absorptionskoeffizienten.
Je nach optischer Weglänge und Konzentration bildet sich ein Verhältnis zwischen dem eintretenden und transmittierten Licht, wie in \prettyref{fig: LBG.png} schematisch dargestellt.

\bild
		{LBG.png}					%Dateiname des Bildes, eindeutiger Name des Bildes, Verweis: \prettyref{fig: Bildname}
		{Schema zur Veranschaulichung des Lambert-Beerschen Gesetzes.%Text, der unter dem Bild erscheint
		\zitat{LBG}}			%Name der Quelle im Literaturverzeichnis
		{scale=0.1}		%angle, draft, scale, height, width
		{h}				%h: here, t: top, b: bottom, p: page; bestimmt, wo das Bild eingefügt wird.

\newpage

Je mehr Licht also absorbiert wird, desto geringer ist die Intensität, welche im Detektor registriert wird.
Dabei ist es wichtig anzumerken, dass die Messung keineswegs eine absolute Angabe über eine Lösung liefert.
Bei einer Messung (in dem Fall mit einem Perkin Elmer Lambda 25 Spektrometer) werden zugleich zwei Proben untersucht,
dabei ist eine Probe eine 'substanzfreie' Lösung, welche die zu untersuchende Substanz nicht beinhaltet und deshalb als Nullabgleich fungiert.
Das liegt an der Tatsache, dass die Küvettenwände und andere Substanzen in der Lösung ebenfalls Licht absorbieren.
In diesem Versuch werden demnach alle gemessenen Absorbanzen von Eisen(II)-Komplexen mit der eisenfreien Lösung verglichen,
indem die Absorbanz des Nullabgleichs -- welcher wie alle anderen Proben Phenanthrolin, Natriumacetat und Ascorbinsäure enthält -- von den gemessenen Absorbanzen abgezogen wird.
Damit werden mögliche Verfälschungen ausgeschlossen und es wird ausschließlich die Absorption der zu vergleichenden Substanz erfasst.\zitat{kp2}

Da freie Eisen(II)-Ionen im sichtbaren Spektralbereich Licht nur geringfügig absorbieren, ist einer direkte photometrische Messung nicht sinnvoll.
Durch eine Komplexierung werden diese in eine farbige Verbindung überführt.
Durch L-Ascorbinsäure als Reduktionsmittel wird sichergestellt, dass Eisen nicht in Form von \ce{Fe^3+} vorliegt.
Damit die Lösung begüglich der pH-Wertes stabil gehalten wird, wird ebenfalls Natriumacetat \ce{NaCH3COO} dazugegeben.
Natriumacetat fungiert dabei als Puffer, sodass kein schwerlösliches Eisen(III)-hydroxid ausfällt.
\cformel[c1]{{Fe^{3+}}_{(aq)} + {3OH^{-}}_{(aq)} -> Fe(OH)3_{(s)}\ce{v}}

Der eigentliche Komplex wird mithilfe von 1,10-Phenanthrolin-Monohydrat (\ce{C12H8N2.H2O}) erzeugt.
Dieses reagiert mit Eisen(II) zu einem Tris(1,10-phenanthrolin)-eisen(II)-Komplex, einem Ferroin-Komplex:
\cformel[c2]{3C12H8N2_{(aq)} + {Fe^{2+}}_{(aq)} -> [Fe(C12H8N2)3]^{2+}_{(aq)}}

Um die Qualität und Verlässlichkeit der photometrischen Messungen zu bewerten, werden die Präzision und die Richtigkeit wie folgt definiert:
Die Richtigkeit beschreibt, wie nah der Mittelwert am tatsächlichen Wert liegt, und wird vor allem durch systematische Fehler (wie ungenaue Stammlösungen oder Gerätedefekte) beeinflusst.
Die Präzision gibt dagegen an, wie nah die einzelnen Messwerte beieinander liegen; hier wirken sich zufällige Fehler wie Pipettiertoleranzen oder Küvettenverschmutzungen aus.
Erst wenn beide Faktoren stimmen, ist die Messung insgesamt genau.\zitat{lol}


\subsubsection{Experimenteller Teil}

\subsubsection{Ergebnisse und Diskussion}


\subsubsection{Zusammenfassung}


\subsubsection{Sicherheit und Entsorgung}
		
Die für den Versuch verwendeten Chemikalien sind in \prettyref{tab: Chemikalien} mit Gefahrensymbolen, H- und P-Sätzen sowie Hinweisen zur Entsorgung aufgeführt.\zitat{test} Verwendete KMR-Stoffe sind in \prettyref{tab: KMR_leer} aufgeführt.\zitat{test}
	
\langetabelle
	{Chemikalien}
	{Verwendete Chemikalien.}
	{H}  % hier ggf. ein großes H eintragen, wenn die Tabelle sonst im Anhang verschwindet
	{l p{1.9cm} p{2.2cm} p{4.5cm} c}
	{Substanz & Symbol & H-Sätze & P-Sätze & Entsorgung \\}
	{ % Erster Stoff

	\ce{NaOH}
		& \begin{tabular}[t]{l} GHS05 \end{tabular}
		& 290, 314
		& 234, 260, 280, \newline 303+361+353, \newline 304+340+310, \newline 305+351+338,
		& (1) \\

	}

\langetabelle
	{KMR_leer}
	{Verwendete kanzerogene, mutagene oder reproduktionstoxische Gefahrstoffe (KMR).}
	{H}
	{l P{3.8cm} P{2.4cm} P{3.5cm} c}
	{Substanz & H-Sätze & KMR-Einstufung & Verwendete Menge & Entsorgung \\}
	{ - & - & - & - & - \\}
	
Entsorgung: \\
\begin{itemize}
	\item[(1)] 
\end{itemize}	


\begin{thebibliography}{4} %manuelles Quellenverzeichnis


	\bibitem{kp}
						{M. Otto}, 
			\textit	{Analytische Chemie}, 
						{4. Auflage}, 
						{Wiley-VCH}, 
						{Weinheim}, 
			\textbf	{2011}, 
			S.			{232-235}.

	\bibitem{kp2}
						{G. Jander, K. Jahr}, 
			\textit	{Maßanalyse}, 
						{17. Auflage}, 
						{Walter de Gruyter}, 
						{Berlin}, 
			\textbf	{209}, 
			S.			{218-219}.

	\bibitem{LBG}
						{Wikiwand}, 
						„{Lay de Beer-Lambert}“, zu finden unter 
			\url		{https://www.wikiwand.com/es/Ley_de_Beer-Lambert}, 
			\textbf	{18.11.2026}.

	\bibitem{lol}
						{Dewetron}, 
						„{Der Unterschied zwischen Präzision, Richtigkeit und Genauigkeit}“, zu finden unter 
			\url		{https://www.dewetron.com/de/news/der-unterschied-zwischen-praezision-richtigkeit-und-genauigkeit/}, 
			\textbf	{18.11.2026}.

	\bibitem{test}
						{DGUV}, 
						„{GESTIS-Stoffdatenbank}“, zu finden unter 
			\url		{https://gestis.dguv.de/}, 
			\textbf	{17.11.2026}.
	
			
\end{thebibliography}

\end{document}

%\newpage

%%%%%%%%%%%%%%%%%%%%%%%%%%%%%%%%%%%%%%%%%%%%%%%%%%%%%%%%%%%%%%%
Neuer Versuch:
\subsection{Versuchsname}

\subsubsection{Einleitung und Aufgabenstellung}
\subsubsection{Theoretischer Hintergrund}
\subsubsection{Experimenteller Teil}
\subsubsection{Ergebnisse und Diskussion}
\subsubsection{Zusammenfassung}
\subsubsection{Sicherheit und Entsorgung}	
Die für den Versuch verwendeten Chemikalien sind in \prettyref{tab: entsorgung1} mit Gefahrensymbolen, H- und P-Sätzen sowie Hinweisen zur Entsorgung aufgeführt.\zitat{Gefahren1} Verwendete KMR-Stoffe sind in \prettyref{tab: KMR1} aufgeführt.
	
\langetabelle
	{entsorgung1}
	{Verwendete Chemikalien.}
	{H}  % hier ggf. ein großes H eintragen, wenn die Tabelle sonst im Anhang verschwindet
	{l p{1.4cm} p{2.2cm} p{4.5cm} c}
	{Substanz & Symbol & H-Sätze & P-Sätze & Entsorgung \\}
	{\\}
    

\langetabelle
	{KMR1}
	{Verwendete kanzerogene, mutagene oder reproduktionstoxische Gefahrstoffe (KMR).}
	{H}
	{l P{3.8cm} P{2.4cm} P{3.5cm} c}
	{Substanz & H-Sätze & KMR-Einstufung & Verwendete Menge & Entsorgung \\}
	{ - & - & - & - & - \\}
Entsorgung: \\
\begin{itemize}
    \item[(1) Die verwendete Lösung wurde mit Salzsäure angesäuert und im Behälter für nicht-nitrathaltige Säuren entsorgt.] 
\end{itemize}	

\newpage

\begin{thebibliography}{4}		%manuelles Quellenverzeichnis
	\bibitem{Gefahren1}
						{H. Heller}, 
			\textit	{Grundlagenpraktikum der Chemie für Studierende der Nanowissenschaften}, 
						{Version 1.18},
                        {Universität Hamburg},
						{Hamburg}, 
			\textbf	{2023}, 
			S.{34}.
			
			
\end{thebibliography} 
\newpage

Nützliche Beispiele:

%Tabelle für Reagenzgläserinhalte
\tabelle
	{Tabellenname}												%Eindeutige Bezeichnung der Tabelle. \prettyref{tab: Tabellenname}
		{Tabellenüberschrift.}	%Inhalt der Reagenzgläser ... bis...
		{htb}	
		{l l l}					
		{Reagenzglas & Inhalt & Beschreibung}				
		{1 & 2~mL NaOH & klare, farblose Lsg.\\
		2 & & 
		}										


%Formel Löslichkeit KClO4 -> chemische Formel
\begin{align}
    \chemtag
    \ce{KClO4(aq) <=> K+(aq) + ClO4-(aq)}
\end{align}

%hier allgemeine Formel 
\formel
		{formel2}				%eindeutige Bezeichnung der Formel, Verweis: \ref{eq: formelname}
		{L_{\mathrm{A}_{a} B_{b}} = \sqrt[m+n]{\frac{K_{\mathrm{L}}}{a^{a} \cdot b^{b}}}  }


Bilder:
	\bild
		{bildname}					%Dateiname des Bildes, eindeutiger Name des Bildes, Verweis: \prettyref{fig: Bildname}
		{Bildbeschreibung.		%Text, der unter dem Bild erscheint
		\zitat{Quelle}}			%Name der Quelle im Literaturverzeichnis
		{Bildeigenschaften}		%angle, draft, scale, height, width
		{Platzierung}				%h: here, t: top, b: bottom, p: page; bestimmt, wo das Bild eingefügt wird. p: spezielle Seite für Abbildungen und Tabellen, Kombinationen möglich, Versuch der Platzierung gemäß Reihenfolge.
	
		
	Bildeigenschaften:
		angle - Bestimmt den Drehwinkel um einen Referenzpunkt für die Grafik. Positive Werte drehen die Grafik gegen den Uhrzeigersinn, negative mit ihm. 
		draft - Ist draft gesetzt, wird nur ein Platzhalter der Grafik angezeigt. Beschleunigt den Übersetzungsvorgang des Dokuments erheblich. 
		scale - Skaliert die Grafik anhand eines Skalierungsfaktors. 
		height - Skaliert die Grafik auf die angegebene Höhe. Die Angabe der Größe muss in einer LaTeX-spezifischen Einheit erfolgen. 
		width - Skaliert die Grafik auf die angegebene Breite. Die Angabe der Größe muss in einer LaTeX-spezifischen Einheit erfolgen. 	
	

Formeln:
	\formel
		{formelname}				%eindeutige Bezeichnung der Formel, Verweis: \ref{eq: formelname}
		{Formel}						%


Tabellen:
	\tabelle
		{Tabellenname}												%Eindeutige Bezeichnung der Tabelle. \prettyref{tab: Tabellenname}
		{Tabellenüberschrift.\zitat{Quelle}}						%Text, der über der Tabelle erscheint, Quelle der Daten
		{Platzierung}													%h: here, t: top, b: bottom, p: page; bestimmt, wo das Bild eingefügt wird. p: spezielle Seite für Abbildungen und Tabellen, Kombinationen möglich, Versuch der Platzierung gemäß Reihenfolge.		
		{Spaltenanzahl_und_Ausrichtung}						%{ l c r }: 3 Spalten, links mittig und rechts; { c c c c }: 4 Spalten, alle mittig
		{Titelzeile}														%Obere Zeile(n), danach folgt ein waagerechter Strich, Spalten durch & trennen, Neue Zeile durch \\
		{Datenzeilen}													%Zeilen unter dem waagerechten Strich, Spalten durch & trennen, Neue Zeile durch \\
		
	Zellen verbinden:	
		\multicolumn{Anzahl_der_verbunden_Spalten}{Ausrichtung (l/c/r}{Inhalt}
	
Reaktionen:
Formeln: 
\begin{center}
\ce{NH+4}, \ce{Fe3(CO)12}, \ce{Al^{3+}},\ce{Cr2O^{2-}7}
\end{center}

Isotope:
\begin{center}
\ce{^{37}_{17}Cl}, \ce{^{32}_{16}S^{2-}}
\end{center}

Bindungen: 
\begin{center}
\ce{H-O-N=S}, \ce{Cl-C#N}
\end{center}

Reaktionsgleichungen: 
\begin{center}
\ce{2H2 + O2 -> 2H2O}
\ce{MgCl2 ->[\text{Elektrolyse}] Mg\,(l) + Cl_{2}\,(g)}
\end{center}

Redoxreaktion:
\begin{center}
$\overset{0}{\ce{S}} + \overset{0}{\ce{O2}} \ce{->} \overset{+\text{IV}}{\ce{S}}+\overset{-\text{II}}{\ce{O2}}$
\end{center}


Literaturverweise:
		\bibitem{zeitschrift}
						{A. Autor}, 
			\textit	{Zeitschrift} 
			\textbf	{Jahr}, 
			\textit	{Band} 
						{Seite-Seite}.
		
		\bibitem{buch_ohne_herausgeber}
						{A. Autor}, 
			\textit	{Titel}, 
						{Verlag}, 
						{Ort}, 
			\textbf	{Jahr}, 
			S.			{Seite-Seite}.
			
		\bibitem{buch_mit_herausgeber}
						{A. Autor}, 
			\textit	{Titel}, 
						{Auflage}, 
						{Verlag}, 
						{Ort}, 
			\textbf	{Jahr}, 
			S.			{Seite-Seite}.
			
		\bibitem{dissertation}
						{A. Autor}, 
						{Dissertation}, 
						{Universität}, 
			\textbf	{Jahr}.
			
		\bibitem{internetseite}
						{A. Autor}, 
						„{Titel}“, zu finden unter 
			\url		{www.xyz.de}, 
			\textbf	{Jahr}.

% Kommentare Sophie:

% <- Für Kommentare oder Code der gerade nicht gebraucht wird, wird beim kompilieren nicht beachtet
%Um Bilder einzufügen: \bild{'Dateiname' (Datei muss im gleichem Ordner mit der TeX Datei liegen)}{'Bildunterschrift'}{width= ...cm (Größe des Bildes)}{H (Ort des Bildes, H=Hier)}

%Für Summenformeln im Text: \ce{z.B.: H2SO4 (Zahlen werden automtisch tief gestellt)}

% Um Formeln im Text zu referieren: \prettyref{fig (für Bilder und Graphen)/ eq (für mathematische Formeln)/ ceq (für chemische Formeln): 'Referenz'}

%Für chemische Formeln: \cformel['Referenz']{2\ox{Fe}{+\RN{3}}Cl3 + \ox{Fe}{+\RN{2}}Cl2 + 8NH4+ + 8 OH- \ce{<=>[{N2H4}]} \ox{Fe}{+\RN{2}}O\cdot \ox{Fe}{+\RN{3}}2O3 \ce{v} + 4 H2O + 8NH4+ + 4 Cl-}

%Für mathematische Formeln: \formel{'Referenz'}{2\cdot d\cdot \ce{sin}(\theta)= n \lambda}

%Tabellen: Online Converter um Exel Tabellen in LaTeX Format zu bringen: https://www.tablesgenerator.com
%\tabelle{'Referenz'}{'Tabellenüberschrift'}{H}{'Breite der Spalte bzw. Textausrichtung' (Spaltenbreite passt sich Textlänge an: l=links, r=rechts, c=zentriert. Spaltenbreite fest und Blocksatz: p{...cm}), Jede Spalte muss einzeln definiert werden, z.B l l l = 3 Spalten}{'Spaltenüberschrift' muss der Anzahl an Spalten entsprechen. Spaltenüberschriften müssen durch ein & getrennt werden. Außerdem wichtig am Ende: \\}{'Tabelleneinträge', müssen ebenfalls durch ein & getrennt werden und am Ende der Zeile ein \\ haben
%	- & - & - & \\
%	- & - & - & \\}
%Beispiel für eine Tabelle (zum Ausführen '%' entfernen):
%\tabelle{Auflösungsvermögen}{Experimentell ermitteltes und theoretisches Auflösungsvermögen für Objektiv 20X und 63X bei schräger und gerader Beleuchtung.}{H}{l l p{2cm} p{3cm}  p{3cm}}{Objektiv&Beleuchtung& Spalt mit einem Peak&experimentelles Auflösungsvermögen~/~nm&theoretisches Auflösungsvermögen~/~nm\\}
%	{20X&gerade&16&>750&1125\\
%	20X&schräg&9&690&563\\
%	63X&gerade&11&>750&600\\
%	63X&schräg&7&584&300\\}